\section{Conclusion}

We tested whether persona-based copy simulation predicts how a real audience ranks
marketing copy, using the Upworthy A/B-test archive as held-out ground truth, and whether the
persona machinery helps at all. Two findings stand out. First, ground-truth reliability is the
binding constraint: most A/B tests lack a statistically distinguishable winner, so validity can
only be assessed on the reliable subset. Second---counter to the premise of persona
simulation---conditioning on personas \emph{degrades} predictive validity: a no-persona zero-shot
ranker achieves a medium-sized rank correlation ($\tau=\baseTau{}$) and \baseTopOne{} top-1
accuracy, while the demographically grounded ten-persona panel reaches only $\tau=\kendallTau{}$
and \topOneAcc{}, with non-overlapping confidence intervals. For predicting aggregate engagement,
the lesson is to query the model's population-level prior directly rather than role-play specific
personas. This does not rule out personas for goals where heterogeneity matters (segmentation,
qualitative objection-finding), but for ranking copy by expected clicks they are, today, worse
than not using them. The artifact-first replication package makes every number reproducible, and
the protocol---reliability filtering, construct-matched elicitation, and a no-persona
baseline---transfers to stronger ground truth and to other copy domains.
